% Part: lambda-calculus
% Chapter: syntax
% Section: eta

\documentclass[../../../include/open-logic-section]{subfiles}

\begin{document}

\olfileid{lam}{syn}{eta}
\olsection{$\eta$-بدلون}

د $\lambda$ ترمونو ترمنځ يوه بله اړيکه هم شته۔ په \olref[fv]{sec} کښې مو
د $\lambd[x][(fx)]$ بېلګه کارولې وه: دا يو آرګومېنټ اخلي او $f$ پرې
تطبيقوي۔ په بله وينا، دا هماغه تابع ده چې~$f$ ده: $\lambd[x][(fx)]N$
او $fN$ دواړه~$fN$ ته راکمېږي۔ د همدې مفکورې د څرګندولو دپاره
$\eta$-راکمول او $\eta$-پراختيا کاروو۔ د برابروالي په دې بېلګه کښې
د تجريد تړلے متغير د تابع په ازادو متغيرونو کښې نۀ وي؛ لاندې تعريف
دا شرط په څرګنده ټاکي۔ \textbf{د سرچينې سپيناوی (OLLAM-033):}
د سرچينې په لومړنۍ وينا کښې دا شرط نۀ دے ويل شوے۔

\begin{defn}[$\eta$-انقباض، $\eredone$]
  \ollabel{defn:beredone}
  \emph{$\eta$-انقباض} ($\eredone$) پر ترمونو تر ټولو کوچنۍ سازګاره
  اړيکه ده چې لاندې شرط پوره کوي:
  \[
    \lambd[x][M x] \eredone M \text{ که } x \notin FV(M)
  \]
\end{defn}

\begin{defn}[$\beta\eta$-راکمول، $\bered$] \ollabel{defn:bered}
  \emph{$\beta\eta$-راکمول} ($\bered$) پر ترمونو تر ټولو کوچنۍ
  انعکاسي او انتقالي اړيکه ده چې $\bredone$ او $\eredone$ پکې
  شاملې دي؛ يعنې د انعکاس او انتقال د قواعدو ترڅنګ لاندې دوه
  قاعدې هم لري:
  \begin{enumerate}
  \item که $M \bredone N$ وي، نو $M \bered N$۔ \ollabel{defn:bered3}
  \item که $M \eredone N$ وي، نو $M \bered N$۔ \ollabel{defn:bered4}
  \end{enumerate}
    
\end{defn}

\begin{defn}
  د $\equal$ معادلتوبي اړيکه د $\eta$-بدلون په دې قاعده غځوو:
  \[
  \lambd[x][f x] \equal f
  \]
  او غځېدلې اړيکه $\equal[\eta]$ ليکو۔
\end{defn}

$\eta$-معادلتوب ځکه مهم دے چې د لامبډا ترمونو د قيمتونو له مخې
برابروالي سره تړاو لري:

\begin{defn}[د قيمتونو له مخې برابري]
  د $\equal$ معادلتوبي اړيکه د (\ext) قاعدې په زياتولو غځوو:
  \begin{center}
  که $Mx \equal Nx$ وي، نو $M \equal N$، په دې شرط چې $x \notin FV(MN)$ وي۔
  \end{center}
  او غځېدلې اړيکه $\equal[\ext]$ ليکو۔
\end{defn}

په لنډه وينا، قاعده وايي چې که دوه ترمونه د تابعو په توګه د يوه او
هماغه آرګومېنټ دپاره يو شان چلند کوي، نو بايد برابر وګڼل شي۔

اوس ثابتوو چې د $\eta$ قاعده د قيمتونو له مخې برابروالي هماغه
اندازه برابروي او تر هغې زيات څۀ نۀ ورزياتوي۔

\begin{thm}
  $M \equal[\ext] N$ هله او يوازې هله چې $M \equal[\eta] N$۔
\end{thm}

\begin{proof}
  لومړے ثابتوو چې $\equal[\eta]$ د قيمتونو له مخې برابروالي د
  قاعدې په وړاندې تړلې ده؛ يعنې د $ext$ قاعده په $\equal[\eta]$
  کښې څۀ نۀ زياتوي۔ له دې امله $\equal[\eta]$، $\equal[\ext]$
  رانغاړي؛ نو که $M \equal[\ext] N$ وي، $M \equal[\eta] N$ هم وي۔
  \textbf{د سرچينې د نښو يادونه (OLLAM-034):} د ثبوت په ځينو
  ځايونو کښې د غځېدلې اړيکې نښه بې‌چپې کرښې ليکل شوې؛ دلته د
  اصلي نسخې رياضيکي نښې ساتل شوې دي او مطلب هماغه قاعده ده۔

  د دې د ثابتولو دپاره چې $\equal[\eta]$ د \ext{} قاعدې په وړاندې
  تړلې ده، پام وکړئ: د هر هغه $M \equal N$ دپاره چې د \ext{} قاعدې
  له لارې ترلاسه شي، $Mx \equal[\eta] Nx$ د هغې مقدمه ده۔ بيا د
  $\equal$ د يوې قاعدې له مخې $\lambd[x][Mx] \equal[\eta]
  \lambd[x][Nx]$ لرو؛ دواړو خواوو ته د $\eta$ په کارولو سره
  $M \equal[\eta] N$ ترلاسه کوو۔

  په همدې ډول ثابتوو چې د $\eta$ قاعده په $\equal[ext]$ کښې
  شامله ده۔ د هر $\lambd[x][Mx]$ او $M$ دپاره چې $x \notin
  FV(M)$ وي، لرو چې $(\lambd[x][Mx])x \equal[\ext] Mx$؛
  له دې څخه د \ext{} قاعدې له مخې $\lambd[x][Mx] \equal[\ext] M$
  ترلاسه کېږي۔
\end{proof}

\end{document}
